\section{Introduction}
Used goods are valued from incomplete descriptions of condition, use and
configuration. Quality uncertainty and voluntary disclosure have a long
economic history~\cite{akerlof1970lemons,grossman1981disclosure,milgrom1981news}.
A predictive price quote introduces a concrete incentive: omitting an
unfavorable recorded attribute can increase the quote. An accurate predictor
on ordinary listings therefore need not remain accurate after sellers choose
which information to disclose. The information available to the platform
matters. Two possible source states of the same listing can produce
identical missing-mileage reports: in one the seller withheld a recorded
value, while in the other the source never recorded that value.
The identifier, other public inputs and peer reports are held fixed.
Figure~\ref{fig:missingness}
shows this boundary. A deployment rule cannot distinguish these cases before
a query. Recorded-value replay may recover a value in the first case and fail
in the second; both attempts consume the shared verification capacity.
This replay is a controlled computational intervention, not physical inspection.

Strategically missing data are already studied in regression and
classification~\cite{zhang2020strategic,krishnaswamy2021withheld}. Costly and
active feature acquisition also have established methods~\cite{janisch2019costly,
shim2018joint,li2021gsm,covert2023information,valancius2024aco}.
The empirical question here is whether learning to suppress concealment-induced
quote increases improves on ordinary error-reduction acquisition and a simpler
robust predictor when listings compete for a fixed batch capacity. Changing one
report can change which other listings receive queries, so a single-listing
information score alone does not describe the deployed interaction.

Strategic classification and causal perspectives model responses to prediction
rules~\cite{hardt2016strategic,dong2018preferences,miller2020causal}; performative
prediction addresses deployment-induced distribution changes~\cite{perdomo2020performative}.
Recent work also jointly designs prediction, audits and
penalties~\cite{ebrahimi2026audit}. Their different actions and outcome models
motivate response analysis but do not establish an incentive property for
continuous quotes with recorded-value replay.
Budgeted randomized audits of strategic reports already minimized incentives
to lie in threshold and top-$k$ allocations~\cite{estornell2021audits}.
Partial verification without payments also has established incentive
characterizations~\cite{ceppi2019verification}.
Estornell et al. check reported types and remove an allocation when a lie
is detected, with an optional fine. Our query recovers one recorded field,
can fail while spending capacity, and changes a continuous quote without
an added penalty. Ceppi et al.'s truthfulness characterization differs from
our finite bound on stored deviation rows. These model distinctions frame
the empirical comparison; they are not new general audit guarantees.

Missing-data mechanisms and predictive treatment of missingness are also
established~\cite{rubin1976missing,lemorvan2021imputation,lemorvan2020neumiss}.
Our indicators cannot identify why a field is blank. Strong categorical
tree models~\cite{prokhorenkova2018catboost} and comparisons with neural tabular
models~\cite{gorishniy2021tabular} motivate retaining a tree control. Acquisition
outputs share the same multi-output ExtraTrees machinery~\cite{geurts2006extratrees},
so an output-specific gain must be assessed against that shared implementation.
Three difficulties remain under this information boundary: natural and
strategic absence are observationally identical, a fitted feedback loop can
reuse outcomes indirectly, and capacity competition couples reports within
a batch. Prediction accuracy and maximum unilateral gain also require
different evaluation endpoints.

We implement two candidates. SAV-Val combines a mask-aware price predictor,
suppression-value acquisition and training on simulated post-query reports.
Its targets are generated through outer folds that contain the entire fitted
feedback loop (Fig.~\ref{fig:overview}; Section~\ref{sec:sav_fitting}).
Counterfactual Deviation-Constrained Mixture Verification (CDM-Val) fixes
the estimator and learns
a distribution over complete capacity-feasible verification rules. It evaluates
signed unilateral payoff differences by reallocating the batch separately for
truthful and withheld reports. Sellers lock their reports before one private
rule is drawn (Section~\ref{sec:cdm_commitment}). Its finite calibration optimum
is distinct from performance on
other listings and peer reports. Randomized response-oracle optimization itself
is established~\cite{jain2011doubleoracle}.

We study public advertised vehicle prices from two countries and one synthetic
source, with stipulated withholding responses. The contribution is a concrete
implementation and an audited assessment of these candidates under a deployable
information boundary. First, we specify public query eligibility and failed
replay with hard capacity (Section~\ref{sec:source_query}). Second, we implement
nested suppression-target fitting and finite mixture calibration
(Sections~\ref{sec:sav_fitting} and~\ref{sec:cdm_commitment}). Third, we evaluate
their distinct inflation and deviation endpoints against matched strong
controls (Section~\ref{sec:results}), retaining unsuccessful comparisons.
It does not introduce strategic learning or auditing as a new
problem, or infer observed seller fraud. The principal finding is heterogeneous
behavior: improvements over a basic masked predictor do not imply improvements
over strong simple controls. All scheduled size/tuning and ordinary-benefit
allowance cases are complete and the local copied evidence passes
reproduction checks. This working draft awaits final author and venue review.

\section{Results}
\label{sec:results}
\subsection{Main comparison and ordinary-benefit reference}
The main nested study is complete: three sources, three fitting seeds and all
six preset capacities, totaling 558 audited evaluation cells. Methods share
1,000 evaluation listings per source. These cohorts were inspected during
historical development, so their retrospective status remains even though the
current fitted loops exclude their labels. Table~\ref{tab:synthetic_development}
reports the prespecified 10\% capacity. MAE denotes mean absolute error in
each source's native price units; errors are not pooled across sources.

On the synthetic source, SAV-Val reduces positive log-quote inflation by
58.19\% relative to masking with ordinary-benefit acquisition. However, it has
54.94\% more inflation than the static robust suppression-output control,
despite 1.28\% lower honest MAE. Its benefit-output feedback counterpart also
has lower mean inflation; the component interval crosses zero.
On JUCars, SAV-Val improves on that feedback counterpart by 14.45\%, but has
34.23\% more inflation than the static robust control. On MUCars, it has
10.92\% less inflation than the static robust control and 1.56\% lower honest
MAE. Its component gain of 8.45\% is below the prespecified 10\% magnitude
criterion. These are distinct comparisons rather than a uniform method gain.

The original ordinary-benefit static robust reference is evaluated separately
in 54 audited cells (Table~\ref{tab:robust_benefit_reference}). SAV-Val has
11.42\% lower inflation than this reference on MUCars, but 92.69\% and
38.65\% higher inflation on synthetic and JUCars. The original fixed-setting
criterion is met only on MUCars. Figure~\ref{fig:capacity} retains every preset
capacity. The synthetic comparison reverses at 30\%, while JUCars favors the
static control at every capacity. MUCars favors SAV-Val even at zero queries,
which points to an estimator-level difference rather than a query-only effect.
These frontier observations do not replace the primary capacity.

\subsection{Finite mixture calibration and source evaluation}
The fixed-estimator CDM study covers all nine source/seed pairs, two price
estimators, six capacities and fourteen rules/mixtures: 1,512 audited cells
and 108 calibration fits. Independent full-matrix linear-program (LP)
solutions agree with the
stored row-generation solutions. All actual rule draws satisfy the hard
capacity. Normal-error constraints do not distinguish rules in this replay
setting because the estimator is fixed and truthful records offer no withheld
recorded value to recover.

At primary capacity, robust-estimator CDM has 18.59\%, 0.73\% and 11.63\%
more positive expected-log inflation than the risk rule on synthetic, JUCars
and MUCars (Table~\ref{tab:cdm_development}). The JUCars interval crosses zero;
the other two favor the risk rule. On MUCars CDM improves on benefit acquisition
by 10.62\%, but has 6.32\% more inflation than uniform mixing. Both price
estimators and all capacities remain visible in Fig.~\ref{fig:cdm_capacity}.
Exact optimality on the stored calibration rows therefore does not establish
a consistent improvement in the held-out mean outcome or universal incentives.

A subsequent common-profile diagnostic exhaustively evaluates legal
unilateral reports on held-out batches under five shared peer-profile
constructions, retaining all 1,512 cells. Rule expectations precede report
maximization. With robust pricing, CDM reduces the upper 5\% envelope mean
relative to the calibration-selected pure rule by 7.34\%, 2.86\% and 6.81\%
on synthetic, JUCars and MUCars, with positive conditional batch intervals.
Those intervals against risk include zero on all three cohorts. Mask-aware
MUCars pricing instead has higher mean and tail envelopes than risk.
This post-hoc diagnostic identifies a bounded mixing effect and leaves the
original primary mean-inflation conclusion unchanged; all references,
capacities and whole-case descriptive tails are retained in the supplement.

\subsection{Selection allowances and comparator boundaries}
The source/size follow-up assigns each method three training-internal
gap-weight trials and twelve residual epochs, then refits the selected model.
All nine source/size cohorts and fixed variants are retained in the
technical supplement: 972 tuning cells and 324 ordinary-benefit reference cells. A separate ordinary-benefit extension selects weights using its
own acquisition rule and constructs its own outer prediction targets and
deployable head. Table~\ref{tab:equal_benefit_reference} includes both fixed and
tuned roles for every independently audited complete three-seed cohort in that
extension. All 27 cases per study and 108 fitting stages per study are audited;
the reference audit includes 54 independent deployment-head refits.

Equal selection allowances do not establish a general full-method gain.
The tuned synthetic comparisons favor the ordinary-benefit static reference
at all three caps. JUCars contrasts change with training size and reference.
Against the ordinary-benefit static reference, tuned JUCars has 20.31\%
and 23.02\% less inflation at caps 1,000 and 2,500, but 27.09\% more
at cap 5,000, with reference-minus-SAV interval $[-0.00333,-0.00125]$.
The honest and attacked error intervals at that largest cap cross zero.

At the configured 5,000 cap (4,774 available training listings), tuned SAV-Val
has 32.71\% more inflation than the tuned static suppression-output control,
with reference-minus-SAV interval $[-0.00368,-0.00157]$. Against its tuned
feedback benefit-output variant, it has 11.83\% less inflation, with interval
$[0.00095,0.00197]$; honest and attacked errors are also lower, with positive
paired error intervals. This component comparison does not establish an
advantage over the static control.

At MUCars cap 1,000, tuned SAV has 7.92\% less inflation than the static
suppression-output control, with a positive conditional interval. Its
1.91\% gain against the feedback benefit-output variant crosses zero,
while honest and attacked errors are larger with intervals excluding zero.

At MUCars cap 2,500, tuned SAV-Val has 19.25\% less inflation than the
static suppression-output control, with reference-minus-SAV interval
$[0.00116,0.00225]$. Honest and attacked MAE are 1.64\% and 2.18\% lower,
with positive paired error intervals. Against the feedback benefit-output
variant, inflation is 16.98\% lower, with interval $[0.00116,0.00171]$;
honest and attacked MAE are also lower. These comparisons retain their
distinct reference identities and do not replace the ordinary-benefit
extension.

Against the tuned ordinary-benefit static reference at MUCars cap 1,000,
SAV-Val has 5.68\% less inflation, with interval $[0.00032,0.00068]$.
This falls below the original 10\% effect-size criterion. Honest and attacked
MAE are 0.23\% and 0.20\% larger, with both error intervals crossing zero.
The fixed-weight comparison has 5.56\% more inflation, with an interval
favoring the reference. Thus weight selection changes this small-cohort
comparison without establishing a general gain.

At MUCars cap 2,500, the tuned ordinary-benefit comparison has 18.92\% lower inflation, with interval $[0.00112,0.00222]$. Honest and attacked MAE are 1.64\% and 2.10\% lower, with positive paired error intervals. This tuned cohort meets the same 10\% inflation-effect and honest-error thresholds. The fixed-weight inflation reduction is 2.61\%, with an interval crossing zero. These results qualify the source/size comparison without establishing a gain across all scheduled cohorts.

At MUCars cap 5,000, tuned SAV-Val reduces inflation by 21.00\% relative to
static suppression-output pricing, with difference 0.001501 and interval
$[0.00101,0.00191]$. Honest and attacked MAE are 1.54\% and 1.83\% lower,
with positive paired intervals. Against the feedback benefit-output variant,
inflation is 10.76\% lower, with interval $[0.00033,0.00103]$; honest and
attacked MAE are 0.60\% and 0.78\% larger, with both error intervals crossing
zero. Thus an inflation improvement does not imply a component accuracy gain.

Against the separately tuned ordinary-benefit static reference at that cap,
inflation is 20.41\% lower, with difference 0.001448 and interval
$[0.00097,0.00186]$. Honest and attacked MAE are 1.54\% and 1.78\% lower,
with positive paired intervals. This cohort meets the prespecified 10\%
inflation-effect and honest-error thresholds. The extension's fixed-weight
comparison has 11.42\% lower inflation, with interval $[0.00045,0.00128]$;
it does not replace the original fixed-setting study's separate interval.
At the same largest cap, tuned inflation is 30.09\% higher on synthetic and
27.09\% higher on JUCars than ordinary-benefit pricing, both with intervals
favoring that reference. All nine cohorts remain visible in the table and
the paired training-size figures; no cross-source success is established.

Selection minimizes honest plus attacked squared transformed-price error,
so tuning can improve accuracy while worsening inflation. No evaluation
outcome selects a weight, seed, source or favorable capacity.

The technical supplement also retains comparator and diagnostic failures.
The acquisition-conditioned oracle control is an explicit one-query regression
adaptation~\cite{valancius2024aco}, not an unchanged classification implementation.
Its strongest declared setting does not establish an inflation advantage for
SAV-Val on the synthetic or Jordanian sources. The full-field SMI adaptation
uses complete categorical groups, declared missing-coordinate domains and a
strict sign-consistency fallback~\cite{zhang2020strategic}. Its large fallback
fractions and errors demonstrate the significance of those assumptions when
extending the numerical formulation; they do not refute the original theorem.

Saved-model sensitivity experiments retain concealment costs, hidden-field
limits and extra natural missingness. A generic disclosure-monotone envelope
eliminates positive inflation but can substantially increase honest error.
Technical missingness strata expose capacity spent on unsuccessful replay;
they are not demographic fairness groups. In the 200-listing JUCars group
with one originally absent optional field, SAV attempts queries on 21.67\%
of listings and fails on 21.50\%, averaged across fitting seeds. This is the
failed-replay branch of Fig.~\ref{fig:missingness} in an observed-record
analysis, rather than evidence of deliberate seller concealment.
Aggregate error can obscure
heterogeneity, as emphasized in group-robust learning~\cite{sagawa2019groups}.
Under common SAV-generated reports,
the static robust control has less inflation in all three sources, although
attacked-error comparisons can favor SAV-Val. Public inference measurements
on a busy shared workstation describe resource use and cannot establish
controlled speedups or physical inspection costs.

\section{Discussion}
The study shows why an improvement over a basic masked predictor is an
insufficient test of strategic valuation. Static robust pricing and ordinary
acquisition remain strong controls. Feedback, acquisition targets and mixture
calibration have effects that depend on the source and capacity. Differences
in approximate responses, fitting samples and target distributions are possible
explanations, not causes established by these experiments.

Inflation is also an incomplete objective. A rule can make withholding
unprofitable by raising honest quotes. The generic envelope control makes this
tradeoff explicit: a zero inflation value need not imply accurate valuation.
Honest and attacked MAE, signed bias, failed queries and remaining deviation
regret must accompany strategic measures. Mixture calibration addresses a
maximum finite deviation criterion, whereas the original primary outcome is
an average positive quote gap; neither should be silently substituted for
the other after observing results.

The sources contain asking prices rather than realized transaction values.
They do not contain observed concealment labels, measured seller utilities or
physical inspection outcomes. Recovered recorded values may themselves be
inaccurate. The stipulated attacker can withhold recorded optional groups,
but cannot fabricate values, change identifiers, choose peers or collude.
Retrospective capped cohorts and ten evaluation batches per source limit
representativeness. Intervals average paired seed contrasts before resampling
the shared batches and condition on the fitted models; they do not quantify
uncertainty over new training populations or real seller behavior.

SAV targets cover the complete outer feedback loop, but teacher and deployment
sample sizes and selected weights can differ. CDM's first source study uses
disjoint estimator, head and probability-calibration rows rather than an outer
replication of the complete mechanism. Its fixed twelve-rule bank and five
calibration peer-profile types are finite restrictions. Prospective independent
markets, measured query costs and observed responses are needed before
deployment. The complete main study supports source-dependent outcomes; the
general empirical increment of either candidate remains unestablished.
